% =====================================================================
%  CHAPTER 4: 青春期发育
% =====================================================================
\chapter{青春期发育}
\setcounter{chapter}{4}
\chapterintro{
掌握青春期基本概念与分期；掌握青春期内分泌调节（HPG轴）及发育机制（难点*）；掌握青春期生长突增的特点和男女差异；掌握青春期性发育过程（Tanner分期）；掌握青春期发育的特殊问题（性早熟、青春期延迟）；熟悉青春期心理行为特点。本章"青春期内分泌"和"发育机制"为教学难点。
}

% ----- 4.1 青春期概述 -----
\section{青春期概述}

\begin{defbox}
\textbf{青春期（Adolescence）}：从青春期发育启动开始至个体生长发育基本停止的时期。WHO定义青春期年龄范围为\keyword{10--19岁}。是人生中生长发育最为剧烈的第二个关键阶段，主要表现为：
\begin{enumerate}[leftmargin=*]
    \item \imp{第二生长突增}——身高、体重快速增长
    \item \imp{性发育}——性器官发育成熟，第二性征出现
    \item \imp{心理行为的巨大变化}
\end{enumerate}
\end{defbox}

\subsection{青春期的分期}

\begin{table}[H]
\centering
\caption{青春期的分期及各期特点}
\begin{tabularx}{\textwidth}{l>{\raggedright\arraybackslash}p{2.5cm}X}
\hline
\textbf{分期} & \textbf{大致年龄范围} & \textbf{主要特点} \\
\hline
\imp{青春早期} & 10--13岁（女10--12岁，男11--13岁）& 生长突增启动；性器官开始发育；第二性征初现 \\
\imp{青春中期} & 13--16岁（女12--15岁，男13--16岁）& 生长突增达高峰（PHV）；第二性征明显；月经初潮/首次遗精 \\
\imp{青春晚期} & 16--19岁 & 生长速度减缓至停止；性发育成熟；生殖功能完善 \\
\hline
\end{tabularx}
\end{table}

\subsection{青春期生长发育的特殊性}

\begin{itemize}[leftmargin=*]
    \item 青春期是全人口中生长发育最为显著的阶段之一
    \item 青春期青少年由于各系统发育尚不成熟，抵抗力较低，容易受到有害因素和不利环境的影响
    \item 在不同生态层面（个人、人际、社区、结构层面），以及在不同发展阶段，青少年的健康都会受到多重生态因素的影响
\end{itemize}

% ----- 4.2 青春期内分泌 -----
\section{青春期内分泌（教学难点*）}

\subsection{下丘脑-垂体-性腺轴（HPG轴）}

\begin{keybox}
\textbf{HPG轴的激活是青春期启动的核心机制：}

\begin{enumerate}[leftmargin=*]
    \item \imp{下丘脑}脉冲式释放\keyword{GnRH（促性腺激素释放激素）}
    \item GnRH刺激\imp{垂体前叶}分泌\keyword{促性腺激素：}
    \begin{itemize}
        \item \textbf{LH（黄体生成素）}
        \item \textbf{FSH（卵泡刺激素）}
    \end{itemize}
    \item LH和FSH作用于\imp{性腺（卵巢/睾丸）}：
    \begin{itemize}
        \item 刺激配子生成（精子/卵子）
        \item 分泌性激素——\textbf{睾酮}（睾丸间质细胞）、\textbf{雌激素}（卵巢颗粒细胞）、\textbf{孕激素}
    \end{itemize}
    \item 性激素通过负反馈调节下丘脑和垂体
\end{enumerate}

\textbf{青春期GnRH分泌模式的变化：}
GnRH的脉冲释放模式在青春期发生显著变化——青春期前GnRH分泌处于抑制状态（低水平、低频脉冲），青春期启动时GnRH脉冲频率和振幅显著增加，这是青春期发动的关键。
\end{keybox}

\subsection{肾上腺皮质功能初现}

\begin{defbox}
\textbf{肾上腺皮质功能初现（Adrenarche）}：青春期前（约6--8岁），肾上腺皮质开始分泌大量雄激素（主要是DHEA和DHEA-S），促进阴毛和腋毛生长，是青春期启动的早期内分泌事件之一，独立于HPG轴的激活。
\end{defbox}

% ----- 4.3 青春期发育机制 -----
\section{青春期发育机制（教学难点*）}

\begin{keybox}
\textbf{青春期启动的调控机制：}

青春期启动的机制尚未完全阐明，目前认为GnRH神经元脉冲释放模式的改变是青春期发动的中枢环节。可能机制包括：

\begin{enumerate}[leftmargin=*]
    \item \imp{中枢抑制的解除}——儿童期存在对GnRH神经元的中枢抑制，青春期启动时该抑制被解除
    \item \imp{兴奋性输入的增加}——kisspeptin（KISS1基因产物）及其受体GPR54对GnRH神经元有强烈的兴奋作用，被认为是青春期启动的关键调控因子
    \item \imp{代谢信号的参与}——瘦素（leptin）可能作为许可信号参与青春期启动，与体脂积累达到一定阈值有关
    \item \imp{表观遗传调控}——DNA甲基化、组蛋白修饰参与青春期启动相关基因的表达调控
    \item \imp{神经递质和神经肽}——谷氨酸（兴奋性）和GABA（抑制性）平衡的改变参与GnRH释放的调控
\end{enumerate}
\end{keybox}

% ----- 4.4 青春期体格发育 -----
\section{青春期体格发育}

\subsection{青春期生长突增}

\begin{keybox}
\textbf{青春期生长突增（Adolescent Growth Spurt）}：青春期的标志性生长现象，也称第二生长突增期。

\textbf{身高突增（Peak Height Velocity, PHV）：}
\begin{itemize}[leftmargin=*]
    \item 身高由每年增长5--10cm开始，逐渐进入突增高峰，一般可达\imp{10--14cm/年}
    \item \imp{男生＞女生}（男生PHV约10--14cm/年，女生约8--12cm/年）
    \item \imp{女生PHV早于男生约2年}（女生约12岁，男生约14岁）
\end{itemize}

\textbf{体重突增：}
\begin{itemize}[leftmargin=*]
    \item 每年增加4--7kg，高峰时可达\imp{8--10kg/年}
    \item 男生以肌肉和骨骼增长为主，女生以脂肪积累为主
\end{itemize}

\textbf{生长停止期：}青春期后期开始，身高基本停止增长，体重也一般停止增长或仅少量增加。
\end{keybox}

\subsection{身体比例与体型变化}

\begin{defbox}
青春期身体比例发生明显变化：
\begin{itemize}[leftmargin=*]
    \item \imp{BMI}：青春期发育使体重快速增加，BMI高百分位数的变化趋势是制定肥胖防控策略的重要依据
    \item 男孩肩宽增加，女孩骨盆增宽
    \item 男性呈倒三角体型，女性呈正三角体型
\end{itemize}

\textbf{Sheldon体型分类法（1940年代）：}
\begin{enumerate}[leftmargin=*]
    \item \imp{内胚层型（Endomorphy）}：身体圆润，消化器官发达
    \item \imp{中胚层型（Mesomorphy）}：身体健壮，骨骼和肌肉发达
    \item \imp{外胚层型（Ectomorphy）}：身体瘦长，神经系统和感觉器官发达
\end{enumerate}
目前更多使用\textbf{Heath-Carter体型分类法}，主要用于运动员体型研究。
\end{defbox}

% ----- 4.5 青春期性发育 -----
\section{青春期性发育}

\subsection{男性性发育}

\begin{keybox}
\textbf{男性性发育顺序（按Tanner分期）：}

\begin{table}[H]
\centering
\caption{男性Tanner性发育分期}
\begin{tabularx}{\textwidth}{l>{\raggedright\arraybackslash}p{3cm}>{\raggedright\arraybackslash}X}
\hline
\textbf{分期} & \textbf{大致年龄} & \textbf{表现} \\
\hline
G1 & 儿童期 & 睾丸体积$<$4mL，阴茎儿童型，无阴毛 \\
G2 & 11--12岁 & 睾丸开始增大（$\geq$4mL），阴囊变薄变红，阴茎开始生长 \\
G3 & 12--13岁 & 睾丸、阴茎继续增大；阴毛增粗、卷曲 \\
G4 & 13--14岁 & 阴茎增粗增长，龟头发育；阴毛接近成人型 \\
G5 & 15--17岁 & 生殖器达成人大小和形态；阴毛呈成人型菱形分布 \\
\hline
\end{tabularx}
\end{table}

\textbf{男性性发育主要事件：}
\begin{enumerate}[leftmargin=*]
    \item \imp{睾丸发育}：最先启动的性发育变化（Tanner G2）
    \item \imp{阴茎发育}
    \item \imp{阴毛发育}
    \item \imp{首次遗精（Spermarche）}：约14--15岁，是男性青春期重要里程碑
    \item \imp{变声、喉结突出}
    \item \imp{面部和体毛生长}
\end{enumerate}
\end{keybox}

\subsection{女性性发育}

\begin{keybox}
\textbf{女性性发育顺序（按Tanner分期）：}

\begin{table}[H]
\centering
\caption{女性Tanner性发育分期}
\begin{tabularx}{\textwidth}{l>{\raggedright\arraybackslash}p{3cm}>{\raggedright\arraybackslash}X}
\hline
\textbf{分期} & \textbf{大致年龄} & \textbf{表现} \\
\hline
B1 & 儿童期 & 仅乳头突出，无乳房组织 \\
B2 & 10--11岁 & 乳房萌芽（breast bud），乳晕增大，为\imp{青春期启动的首个体征} \\
B3 & 11--12岁 & 乳房和乳晕进一步增大 \\
B4 & 12--13岁 & 乳晕和乳头突出形成第二小丘 \\
B5 & 14--15岁 & 成熟乳房，乳晕与乳房轮廓融合 \\
\hline
\end{tabularx}
\end{table}

\textbf{女性性发育主要事件：}
\begin{enumerate}[leftmargin=*]
    \item \imp{乳房发育}：青春期启动的最早体征（Tanner B2），约10--11岁
    \item \imp{阴毛发育}：略晚于乳房发育
    \item \imp{月经初潮（Menarche）}：约12--13岁，是女性青春期最重要的里程碑事件
    \item \imp{骨盆增宽}、皮下脂肪增厚（形成女性体型）
\end{enumerate}

\textbf{月经初潮的影响因素：}遗传（最主要）、营养状况、体重/体脂含量、社会经济条件、地理气候、体育锻炼强度等。
\end{keybox}

\subsection{性发育的影响因素}

\begin{defbox}
性发育受多种因素综合影响：
\begin{enumerate}[leftmargin=*]
    \item \imp{遗传因素}——决定青春期启动时机的最重要因素
    \item \imp{营养状况}——营养过剩可导致青春期提前
    \item \imp{体脂含量}——瘦素是脂肪组织分泌的激素，达到一定水平是青春期启动的许可信号
    \item \imp{环境内分泌干扰物（EEDs）}——可干扰正常性发育进程，导致性早熟或延迟
    \item \imp{社会心理因素}——家庭压力、社会应激等
    \item \imp{体育锻炼}——高强度训练可能导致青春期延迟
\end{enumerate}
\end{defbox}

% ----- 4.6 青春期特殊问题 -----
\section{青春期发育的特殊问题}

\subsection{性早熟}

\begin{keybox}
\textbf{性早熟（Precocious Puberty）}：女孩8岁前、男孩9岁前出现第二性征发育。

\begin{itemize}[leftmargin=*]
    \item \imp{真性性早熟（中枢性）}：HPG轴提前激活，GnRH脉冲释放过早启动。可特发性或继发于中枢神经系统病变
    \item \imp{假性性早熟（外周性）}：性激素来源于性腺或肾上腺的自主分泌（如肿瘤），或外源性激素暴露（EEDs），HPG轴未激活
    \item \imp{不完全性性早熟}：单纯乳房早发育、单纯阴毛早现、单纯月经早现
\end{itemize}

\textbf{危害：}
\begin{enumerate}[leftmargin=*]
    \item 骨骺早期闭合，最终身高偏矮
    \item 心理社会适应问题（与同龄人发育不同步）
    \item 可能的潜在病因（如肿瘤）
\end{enumerate}
\end{keybox}

\subsection{青春期延迟}

\begin{defbox}
\textbf{青春期延迟（Delayed Puberty）}：女孩13岁、男孩14岁仍无青春期启动的任何体征。

常见原因：
\begin{enumerate}[leftmargin=*]
    \item \imp{体质性青春期延迟}（最常见）——有家族史，最终发育正常
    \item \imp{低促性腺激素性性腺功能减退}——GnRH/LH/FSH分泌不足
    \item \imp{高促性腺激素性性腺功能减退}——性腺本身功能障碍（如Turner综合征、Klinefelter综合征）
    \item 慢性疾病、营养不良、过度运动等
\end{enumerate}
\end{defbox}

\subsection{青春期心理卫生}

\begin{defbox}
青春期是心理行为发展的关键期，常见心理卫生问题包括：
\begin{enumerate}[leftmargin=*]
    \item 自我认同危机
    \item 情绪波动（与性激素变化相关）
    \item 逆反心理
    \item 同伴压力
    \item 冒险行为倾向
    \item 体象关注（对身体外貌过度关注）
\end{enumerate}
\end{defbox}

% ----- 复习题 -----
\section{复习题}

\begin{enumerate}[leftmargin=*, itemsep=8pt]
    \item \textbf{【名词解释】}青春期；第二生长突增；PHV（身高突增高峰）
    \item \textbf{【名词解释】}月经初潮；性早熟；青春期延迟
    \item \textbf{【名词解释】}肾上腺皮质功能初现；GnRH
    \item \textbf{【名词解释】}Tanner分期；HPG轴
    \item \textbf{【简答题】}简述青春期的分期及各期特点。
    \item \textbf{【简答题】}试述下丘脑-垂体-性腺轴（HPG轴）的调控机制及在青春期启动中的作用。
    \item \textbf{【简答题】}简述青春期生长突增的特点及男女差异。
    \item \textbf{【简答题】}简述男性Tanner性发育分期的要点。
    \item \textbf{【简答题】}简述女性Tanner性发育分期的要点。
    \item \textbf{【简答题】}试述性早熟的定义、分类及主要危害。
    \item \textbf{【简答题】}简述青春期发育的影响因素。
\end{enumerate}

\subsection*{参考答案要点}

\begin{enumerate}[leftmargin=*, itemsep=5pt]
    \item \textbf{青春期}：从发育启动至生长发育基本停止（10--19岁，WHO）。\textbf{第二生长突增}：青春期身高体重快速增长的标志性现象。\textbf{PHV}：身高增长速率达到峰值的年龄，女生约12岁、男生约14岁。

    \item \textbf{月经初潮}：女性第一次月经来潮，约12--13岁。\textbf{性早熟}：女$<$8岁、男$<$9岁出现第二性征。\textbf{青春期延迟}：女$>$13岁、男$>$14岁无青春期启动体征。

    \item \textbf{肾上腺皮质功能初现}：青春期前肾上腺开始分泌大量雄激素，促进阴毛腋毛生长，是青春期启动早期事件之一。\textbf{GnRH}：下丘脑分泌的促性腺激素释放激素，脉冲式释放，是青春期启动的关键激素。

    \item \textbf{Tanner分期}：由Tanner提出的性发育分期标准，分5期。\textbf{HPG轴}：下丘脑-垂体-性腺轴，是调控青春期发育和生殖功能的核心内分泌轴。

    \item 三期：青春早期（生长突增启动、性发育初现）、青春中期（PHV高峰、性发育明显、初潮/遗精）、青春晚期（发育趋缓至停止、性成熟）。

    \item GnRH脉冲释放→垂体分泌LH/FSH→性腺分泌性激素。青春期启动核心：GnRH脉冲频率和振幅增加。中枢抑制的解除和kisspeptin/GPR54系统的激活是关键。

    \item 身高由年增5--10cm至PHV达10--14cm/年，男生＞女生；女生PHV早于男生约2年。体重每年4--7kg，峰时可达8--10kg/年。

    \item Tanner G1--G5：G1儿童型→G2睾丸开始增大（$\geq$4mL）→G3阴茎睾丸继续增大，阴毛增粗→G4龟头发育→G5成人型。首次遗精约14--15岁。

    \item Tanner B1--B5：B1仅乳头突出→B2乳房萌芽（最早体征）→B3乳房乳晕增大→B4乳晕乳头形成第二小丘→B5成熟乳房。月经初潮约12--13岁。

    \item 定义：女$<$8岁、男$<$9岁出现第二性征。分类：真性（中枢性，HPG轴提前激活）、假性（外周性，性激素来源异常）、不完全性。危害：骨骺早闭致矮小、心理问题、可能潜在疾病。

    \item 遗传（最主要）、营养状况、体脂含量/瘦素、环境内分泌干扰物（EEDs）、社会心理因素、体育锻炼强度等。
\end{enumerate}

\newpage
